\section{Bubbling analysis}\label{sec:bubbling}
        	 	
        	 	
        	 	In this section, we prove Theorem 5 and in Section 4 we complete the proof of Theorem 1. Before we begin either of these proofs notice that given $y\in {\jepPubScript{Y}}\smallsetminus \partial N$ the bubbling and neck analysis in \cite{BS17} goes through precisely as before since it is of a local nature. In particular the conclusions of Theorems 1 and 5 hold over smooth domains $U\cemb N\smallsetminus\partial N$ such that ${\jepPubScript{Y}}\cap \partial U =\varnothing$ and $y\in {\jepPubScript{Y}}\smallsetminus \partial N$ respectively, which implies that the proofs of both of these theorems is complete in the case that ${\jepPubScript{Y}}\cap \partial N$ is empty. We need only focus on the setting $y\in {\jepPubScript{Y}}\cap \partial N$ from now on. 
        	 	
        	 	Our main result here is a blow-up theorem (Theorem \ref{thm:blowup} below) which will detect a non-trivial bubble or half-bubble in all regions of coalescing index at $\partial N$. 
        	 	This result is based on a localized version of the compactness result for free boundary minimal hypersurfaces from \cite{ACS17}. Once we have proved Theorem \ref{thm:blowup}, we can then follow a scheme related to the bubbling analysis for closed minimal hypersurfaces developed in \cite[\S3]{BS17} to construct various point-scale sequences which will detect \emph{all} the non-trivial bubbles and half-bubbles that develop at points of curvature concentration on $\partial N$, yielding Theorem 5. 
        	 	
        	 	In order to obtain the quantization result in Theorem 1 from our proof of Theorem 5, it remains to show that no curvature is lost in \emph{annular neck regions} between the bubble scales. This last step will be carried out in Section 4.\\
        	 	
        	 	
        	 	
        	 	In order to state our results in a unified manner, we will use the following notation and conventions throughout. Consistently with Section \ref{sec:prelim} denote by $\Pi_1(a):=\{x^1\geqslant a\}$ a closed half-space of $\R^{n+1}$ \emph{containing the origin} and by $\Pi(a)$ its boundary (though we drop the dependence on $a$ when it is irrelevant).  We will consider %
        	 	a relatively open sub-domain $C^{\R^{n+1}}_\varrho(a)\subset \Pi_1$,  %
        	 	defined by $C^{\R^{n+1}}_\varrho(a)=[a,a+\varrho)\times B^{\R^n}_\varrho (0)$. We will allow $\varrho=\infty$ in the sequel at which point the domain is simply $\Pi_1(a)$.  %
        	 	Now equip $C^{\R^{n+1}}_\varrho$ with a smooth metric $g$ such that it is Euclidean at $(a,0,\dots,0)$ and for all $z\in C^{\R^{n+1}}_\varrho\cap \Pi$, $\partial_{x^1}(z)$ is the inward unit normal to $\Pi$ with respect to $g$. We say that $(C^{\R^{n+1}}_\varrho,g)$ is an \emph{adapted} domain and in the sequel we will allow the parameters $\varrho$ and $g$ to vary under these constraints, and $a$ will largely remain fixed.   %
        	 	
        	 	
        	 	Such domains $(C^{\R^{n+1}}_\varrho,g)$ can be used to analyze local properties of free boundary minimal hypersurfaces in arbitrary $N$ without loss of generality via the use of
        	 	a \emph{Fermi-coordinate} neighborhood at the boundary of $N$: 
        	 	Let $-\nu$ be the inward pointing normal to $\partial N$ and $\gamma_1(N)>0$ be so that 
        	 	\begin{eqnarray*}
        	 		& F:[0,\gamma)\times \partial N \longrightarrow N&\\
        	 		&(t,z)\longmapsto \operatorname{Exp}^N_z (-t\nu)&
        	 	\end{eqnarray*} 
        	 	is injective for $t\in [0,\gamma_1)$. 
        	 	
        	 	Consider a normal coordinate neighborhood of $q$ in $\partial N$, $\{x_i\}_{i=2}^{n+1}$. For ease of notation, we will now choose $\gamma_0(N) >0$ to be smaller than both $\gamma_1$ defined above, and the injectivity radius of $\partial N$ so that, first of all, these boundary normal coordinates are defined on $B^{\partial N}_\gamma (q)$ for any $q\in \partial N$ and $\gamma \leqslant \gamma_0$. Second of all, for all $q\in \partial N$ and $\gamma\leqslant \gamma_0$, we may now extend these to a Fermi-coordinate neighborhood $C_{\gamma}(q) \cong [a,a+\gamma)\times B^{\R^n}_\gamma (0)$ via the exponential map giving $x^{1} = a+ \operatorname{Exp}^N_{(0,x^2,\dots,x^{n+1})} (-t\nu)$ for $0\leqslant t<\gamma$. 
        	 	
        	 	Notice that the metric $g$ in these coordinates coincides with the Euclidean metric at $(a,0,\dots,0)$. Moreover $\partial_{x^1}(z) = -\nu(z)$, in particular $g_{1j}(z)=0$, for all $z\in \partial N = \{x^1=a\}$ and $j\geqslant 2$. In particular $C_\gamma(q)\cong C^{\R^{n+1}}_\gamma(a)$ is an adapted domain.  \\
        	 	
        	 	For $V\subset (C^{\R^{n+1}}_\varrho,g)$ consider the set $\mathfrak{M}^{V}$ (which shall depend on both $C^{\R^{n+1}}_\varrho$ and the background metric $g$) of smooth, connected and properly embedded minimal hypersurfaces $P\subset V$, furthermore requiring that if $P\cap \Pi \neq \varnothing$ then $P$ is free boundary with respect to $\Pi$. At the level of regularity, we always tacitly assume $V$ to be the intersection of a \emph{smooth} domain in $\R^{n+1}$ with $\Pi_1$. Any variational properties of $P$ are computed with respect to compactly supported variations in $V$ -- i.e. free boundary variations on $V\cap \Pi$ as well as Dirichlet conditions on $\partial V\smallsetminus \Pi$. Following our introduction, we define
        	 	\begin{equation*}
        	 	\mathfrak{M}^{V}_p(\Lambda,\mu) = \big\{ P \in \mathfrak{M}^{V} :  \forall x\in C^{\R^{n+1}}_\varrho, R>0,  \ \h^n(P\cap B^{\R^{n+1}}_R(x))\leqslant \Lambda R^n, \ \text{ and }\lambda_p(P) \geqslant -\mu \big\}.
        	 	\end{equation*}
        	 	
        	 	
        	 	
        	 
        	 	We recall from \cite[Th.~29]{ACS17} that if $2\leqslant n\leqslant 6$ and $\{M_k\}\subset \mathfrak{M}_p(\Lambda, \mu) \subset \mathfrak{M}$ then there exists a smooth, connected, compact embedded minimal hypersurface $M\subset N$ meeting $\partial N$ orthogonally along $\partial M$, $m\in \mathbb{N}$ and a finite set ${\jepPubScript{Y}}\subset M$ with cardinality $|{\jepPubScript{Y}}|\leqslant p-1$  such that, up to subsequence, $M_k \to M$ locally smoothly and graphically on $M\smallsetminus {\jepPubScript{Y}}$ with multiplicity $m$. To avoid ambiguities, let us remark that from now on we will always assume that  ${\jepPubScript{Y}}$ is the \emph{minimal} such set, so that in particular the convergence is never smooth about any $y\in{\jepPubScript{Y}}$.
        	 	
        	 	We will require the following local version of that compactness theorem.
        	 	
        	 	\begin{lem}\label{lemma:local}
        	 		Let $2\leqslant n \leqslant 6$ and $\{(C_{k}, g_k)\}$ be a sequence of adapted domains where we set $C_k=C^{\R^{n+1}}_{\varrho_k}(a)$ with $\{\varrho_k\}$ monotonically increasing, and choose $\eps >0$ small enough that $B^{\R^{n+1}}_\eps(0)\cemb C_1$. We will assume that $g_k\to g$ smoothly for some limit metric $g$ on any relatively open subset $V\cemb C_k$. Suppose we have a sequence $P_k\in \mathfrak{M}^{C_k}_p(\Lambda,\mu)$ (for some fixed constants $\Lambda, \mu\in \R_{>0}$ and a positive integer $p$ independent of $k$) such that $P_k\cap B^{\R^{n+1}}_{\eps}(0)\neq \varnothing$.
        	 		Then, up to subsequence: 
        	 		
        	 		
        	 		\begin{enumerate}
        	 			\item For any relatively open $V$ such that $B_\eps^{\R^{n+1}}(0)\subset V\cemb C_k$ for some $k$, there exists a smooth, connected, and embedded minimal hypersurface $P\subset V$ where $P_k \to P$ locally smoothly and graphically, with multiplicity $m\in \N$, for all $x\in P\smallsetminus {\jepPubScript{Y}}$ where ${\jepPubScript{Y}}\subset P$ is a finite set with cardinality $|{\jepPubScript{Y}}|\leqslant p-1$. 
        	 			
        	 			 If $\partial P\neq\varnothing$, then $P$ meets $\Pi$ orthogonally along $\partial P$ and if $P \in \mathfrak{M}^V$, then $P \in \mathfrak{M}^V_p(\Lambda,\mu)$.
        	 			 
        	 			  Finally ${\jepPubScript{Y}}\neq \varnothing$ if and only if there exists $\{y_k\}\subset V$ with $y_k\to y\in P$, and $r_k\to 0$ so that $\operatorname{index}(P_k\cap B^{\R^{n+1}}_{r_k}(y_k))\geqslant 1$. In either case we say that $y\in {\jepPubScript{Y}}$. 
        	 			
        	 			
        	 			
        	 			\item Assuming that $\varrho_k \to \infty$, then there exists a limit $M$ which is a smooth, embedded  minimal hypersurface in $(\Pi_1,g)$.   
        	 		If we additionally assume $\liminf_{k\to\infty}\lambda_p(P_k)\geqslant 0$ then  $P \in \mathfrak{M}^{\Pi_1}$ implies $P \in \mathfrak{M}^{\Pi_1}(\Lambda,p-1)$.
        	 			
        	 			Further assuming $g=g_0$ is the Euclidean metric then $P$ is either a bubble or a half-bubble; if $P$ is not properly embedded, then $P\equiv \Pi$. 
        	 			
        	 			Finally if ${\jepPubScript{Y}}\neq \varnothing$ then $P$ must be a plane in $\Pi_1$ or a half-plane orthogonal to $\Pi$. 
        	 			
        	 		\end{enumerate}
        	 		
        	 	\end{lem}
        	 	
        	 	\begin{proof}
        	 		Part (1):        		 All the steps in the proof of \cite[Th.~29]{ACS17} are local, hence can be adopted almost verbatim with only some cosmetic changes to conclude all but the final statement concerning the equivalent characterization of the condition ${\jepPubScript{Y}}\neq\varnothing$. To prove that assertion, let us first observe that since $P$ is smooth and $V$ is compact we can pick $r>0$ so that $P\cap B^{\R^{n+1}}_r(z)\cap V$ is \emph{strictly} stable for all $z\in P$. Thus, if ${\jepPubScript{Y}}=\varnothing$ then the convergence would be smooth everywhere, hence in particular for $k$ large enough each hypersurface $P_k$ would be stable in all balls of radius $r/2$ centered at points in $V$ (as defined above) by smooth convergence. Hence no sequence as in the statement can actually exist. Instead, for the converse implication one just needs to invoke the interior and free boundary curvature estimates of Schoen-Simon \cite{SS81} (cf. \cite[Th.~19]{ACS17}).  
        	 		
        	 		\
        	 		
        	 		Part (2):
        	 		We can apply Part (1) to some compact exhaustion $\{V_{\ell}\}$ of $\Pi_1$ to conclude the first statement via a diagonal sequence argument.%
        	 		
        	 		If $P$ is proper, i.e. $P \in \mathfrak{M}^{\Pi_1}$, then by Part (1) $P\in \mathfrak{M}^V_p(\Lambda, \mu)$ for \emph{all} $V \subset \Pi_1$. We now check that in this case in fact $\operatorname{index}(P)\leqslant p-1$. Assuming the contrary guarantees the existence of a set $V$ with $\lambda_p(P\cap V)=\alpha<0$ and therefore, following the argument in \cite[p.~2599--2600]{ACS15}, we find that eventually $\lambda_p(P_k\cap V)\leqslant \alpha/2$, contradicting the assumption that $\liminf_{k\to\infty}\lambda_p(P_k)\geqslant 0$.  
        	 		
        	 		If $g=g_0$ is the Euclidean metric, then we distinguish two cases. If $P$ is proper, then exploiting the assumption of Euclidean volume growth, we conclude by Proposition \ref{pro:equiv} that it has finite total curvature. If on the other hand $P$ is not proper, i.e. an interior point of $P$ lies in $\Pi$, then by a maximum principle argument $P=\Pi$, and therefore $P$ trivially has finite total curvature. 
        	 		
        	 		Finally, if ${\jepPubScript{Y}}\neq \varnothing$ then we may as well assume that $P$ is properly embedded (if not then $P=\Pi$ and we are done). Since $P$ is properly embedded and $\Pi_1$ is simply connected $P$ is two-sided. When we couple this with the lack of smooth convergence, we conclude that the convergence must be of multiplicity $m\geqslant 2$ (via Allard's interior and boundary regularity cf. \cite[Ths.~5, 17]{ACS17}) which allows for the construction of a positive Jacobi field over compact subsets $W\cemb \Pi_1$ (not necessarily vanishing at the boundary of $W$). In turn, this implies that $\lambda_1(P\cap W)\geqslant 0$ and the limit is stable over all such $W$ (here we use the standard argument to establish the positiveness of a first eigenfunction, noting that both an interior and boundary Hopf maximum principle will be required); hence $P$ is stable in $\Pi_1$. The conclusion follows by the usual classification of stable hypersurfaces with Euclidean volume growth when $P$ has no boundary \cite{SS81}, and Remark \ref{rmk:higher-stable} when $P$ has a free boundary on $\Pi$. 
        	 		\end{proof}
        	 	
        	 	
        	 	We can now establish the first main result of this section as a consequence of the above statement as well as the localized compactness theory developed in \cite[\S2]{BS17} for the case of closed hypersurfaces. 
        	 	
        	 	 We have already seen in Lemma \ref{lemma:local} that bad points of convergence $y\in {\jepPubScript{Y}}$ correspond to coalescing regions of index, so these are the points that we analyze in detail here. We will denote by $B_r(p)$ a normal coordinate neighborhood of $p\in N$ and by $C_r(q)$ a Fermi-coordinate neighborhood of $q\in \partial N$. 
        	 
